% !TeX spellcheck = en_US
\documentclass[french]{yLectureNote}

\title{Électromagnétisme 1 PS}
\subtitle{Physique}
\author{Paulhenry Saux}
\date{\today}
\yLanguage{Français}
%lionels.calmels@cemes.fr
\professor{F.Pettinari}
\usepackage{graphicx}%----pour mettre des images
\usepackage[utf8]{inputenc}%---encodage
\usepackage{geometry}%---pour modifier les tailles et mettre a4paper
%\usepackage{awesomebox}%---pour les boites d'exercices, de pbq et de croquis ---d\'esactiv\'e pour les TP de PC
\usepackage{tikz}%---pour deiffner + d\'ependance de chemfig
% \usepackage{tabularx}%---pour dimensionner automatiquement les tableaux avec variable X
\usepackage{awesomebox}%---Pour les boites info, danger et autres
\usepackage{menukeys}%---Pour deiffner les touches de Calculatrice
\usepackage{fancyhdr}%---pour les en-t\^ete personnalis\'ees
\usepackage{blindtext}%---pour les liens
\usepackage{hyperref}%---pour les liens (\`a mettre en dernier)
\usepackage{caption}%---pour la francisation de la l\'egende table vers Tableau
\usepackage{pifont}
\usepackage{array}%---pour les tableaux
\usepackage{lipsum}
\usepackage{yFlatTable}
\usepackage{multicol}
\newcommand{\Lim}[1]{\lim\limits_{\substack{#1}}\:}
\renewcommand{\vec}{\overrightarrow}
\newcommand{\N}[0]{\mathbb{N}}
\newcommand{\dd}{\mathrm{d}}
\newcommand{\norm}[1]{||\vec{#1}||}
\newcommand{\fo}{\psi(\vec{r},t)}
\newcommand{\foe}{\psi(\vec{r},t)\*}
\newcommand{\HH}{\hat{H}}
\newcommand{\hb}{\hbar}
\newcommand{\lap}{\nabla^2}
\newcommand{\lapcc}{\frac{\partial^2 }{\partial x^2}+\frac{\partial^2 }{\partial y^2}+\frac{\partial^2 }{\partial z^2}}
\newcommand{\E}{\vec{E}(M)}
\begin{document}
\chapter{Champ électrique et potentiel par des charges ponctuelles}
\section{Charges ponctuelles}
\begin{definition}[Charge ponctuelle]
Les charges ponctuelles sont des particules élémentaires, comme l'électron \(e^-\) avec \(q_e = 1.6\cdot 10^{-19}C, q_p = + 1.6\cdot 10^{-19}C\)
\end{definition}
\section{Lois de Coulombs}
Elles décrivent les forces que les charges exercent les unes sur les autres. On note q la charge au point M et \(q_p\) la charge au point M.
\begin{theorem}[Lois de Coulombs]
La force que la charge \(q_p\) exerce sur q est \[\vec{F_{q_p\to q}} = \frac{qq_p}{4\pi \epsilon_0}\frac{\vec{e_{PM}}}{\norm{PM}^2}\] avec \(\epsilon_0\) la permittivité du vide \(\frac{1}{36\pi \cdot 10^{9}}\) et \(\frac{1}{4\pi \epsilon_0}\simeq 9\cdot 10^{9}\)
\end{theorem}
%TODO Schema 1

Si 2 charges ponctuelles sont de m\^eme signe, elles se repoussent, dans le cas contraire, elles s'attirent. On a les m\^emes propriétés qu'en mécanique du point.
\section{Champ électrique crée par une charge ponctuelle}
\subsection{Définition}
\begin{definition}[Champ électrique]
\[\vec{F_{q_p\to q}} = q \cdot \E\] avec \(\E\) qui ne dépend que de la charge \(q_p\) : \[\E = \frac{q_p}{4\pi\epsilon_0}\frac{\vec{e_{PM}}}{\norm{PM}^2}\] est le champ électrique crée par \(q_p\) au point M
\end{definition}
\subsection{Remarques}
\begin{itemize}
 \item \([\E] = \frac{[charge]}{[\epsilon_0][L^2]}\)
 \item \(\norm{E} \sim \frac{1}{\norm{PM}^2}\)
 \item Le champ \(\E\) fuit les charges positives et est orientée vers les charges négatives. La direction des vecteurs du champ est radiale.
 \item Unité du champ électrique : \(NC^{-1} / Vm^{-1}\)
\end{itemize}
\subsection{Théorème de superposition}
\begin{theorem}[Théorème de superposition]
Le champ électrique d'une distribution de charges ponctuelles \(q_i\)  vaut \(\sum_{i=1}^n \vec{E_i}(M)\) au point M.
\end{theorem}
\section{Potentiel électrique}
\subsection{Circulation de \(\E\)}
Soit q placée en O qui crée en M \(\E\) : \[\E = \frac{q}{4\pi \epsilon_0 r^2}\vec{e_r}\]

Calculons la Circulation entre A et B (on reste en coordonnée sphérique pour des raisons de symétrie) :
\begin{flalign*}
\mathcal{C}_{A\to B}\E &= \int_A^B\E\cdot \dd \vec{l}(M)\\
&= \int_A^B\E\cdot (\dd r \vec{e_r}+r\dd \theta \vec{e_{\theta}}+r\sin(\theta)\dd \varphi \vec{e_{\varphi}})\\
&= \int_A^B\frac{q}{4\pi \epsilon_0 r^2}\dd r\\
&= \frac{q}{4\pi \epsilon_0}\int_A^B\frac{\dd r}{r^2}\\
&= \frac{q}{4\pi \epsilon_0}\int_{r_A}^{r_B}\frac{\dd r}{r^2}\\
&= \frac{q}{4\pi \epsilon_0 r_a} - \frac{q}{4\pi \epsilon_0 r_b}
\end{flalign*}
Le résultat ne dépend que de la position des points et non du chemin suivi, le champ est donc conservatif
\subsection{Potentiel électrique}
On peut écrire \(\mathcal{C}_{A\to B}\E = V(A)-V(B)\) avec \(V(M) = \frac{q}{4\pi\epsilon_0 r} + Cst\). V(M) est le potentiel électrique crée au point M par la charge ponctuelle q placée au point O.
\subsection{Remarques}
\begin{itemize}
 \item \(V(M) \sim \frac{1}{\norm{PM}}\)
 \item \([V(M)] = \frac{[charge]}{[\epsilon_0]L}\)
 \item \(V(M)\) est défini à une constante additive arbitraire\marginInfo{C'est pas grave car ce qui a un sens physique, c'est la différence de potentiel où la constante disparaît.}.
\end{itemize}
\subsection{Théorème de superposition}
De la m\^eme façon que pour la force, le potentiel obéit aux m\^emes lois.
 \end{document}
